\chapter{Sécurité et règles d'accès (RLS)}

\section{Modèle de sécurité}
La sécurité repose sur trois piliers~:
\begin{enumerate}
  \item \textbf{Authentification JWT} via Supabase Auth.
  \item \textbf{Rôles séparés} dans \texttt{user\_roles}, vérifiés par des
        fonctions \texttt{SECURITY DEFINER} (\texttt{has\_role},
        \texttt{is\_parent\_of}) --- évite l'élévation de privilèges et la
        récursion RLS.
  \item \textbf{Row-Level Security (RLS)} activée sur \emph{toutes} les tables du
        schéma \texttt{public}, complétée par des \texttt{GRANT} explicites.
\end{enumerate}

\section{Principes des politiques RLS}
\begin{longtable}{p{4.5cm}p{10cm}}
\toprule \textbf{Domaine} & \textbf{Règle d'accès} \\ \midrule \endhead
Données personnelles & Chaque utilisateur ne voit/modifie que ses propres lignes
(\texttt{auth.uid() = user\_id}) : \texttt{student\_scores},
\texttt{ai\_generated\_content}, \texttt{ai\_lesson\_comments},
\texttt{chat\_conversations}, \texttt{chat\_usage}. \\
Contenu pédagogique & Lecture réservée aux utilisateurs authentifiés ; gestion
(écriture) réservée aux \texttt{pedago} et \texttt{admin} : \texttt{chapters},
\texttt{lessons}, \texttt{chapter\_exercises}, \texttt{chapter\_quizzes},
\texttt{exams}, \texttt{filieres}. \\
Suivi parental & Un parent accède aux données de ses enfants \emph{uniquement}
si une liaison active existe (\texttt{is\_parent\_of}) :
\texttt{parent\_reports}, \texttt{ai\_lesson\_comments} (enfants). \\
Administration & Les administrateurs disposent de politiques \og manage\fg{}
(ALL) sur les codes, abonnements, journaux, etc. \\
Audit & \texttt{activity\_logs} : insertion par utilisateurs authentifiés,
lecture par le propriétaire et les administrateurs. \\
\bottomrule
\end{longtable}

\section{Bonnes pratiques appliquées}
\begin{itemize}
  \item \textbf{Aucun secret en base} : les clés d'API vivent dans les secrets
        Supabase, jamais dans les tables ni dans le client.
  \item \textbf{\texttt{service\_role} jamais exposé} au navigateur ; utilisé
        seulement à l'intérieur des Edge Functions.
  \item \textbf{Validation des réponses côté serveur} via des fonctions
        \texttt{SECURITY DEFINER} (\texttt{check\_exercise\_answer},
        \texttt{check\_quiz\_answer}) --- l'élève ne peut pas connaître la réponse
        en inspectant le client.
  \item \textbf{GRANT explicites} pour chaque table conformément aux exigences de
        PostgREST.
\end{itemize}

\begin{infobox}[\faExclamationTriangle\ Point de vigilance]
Les rôles ne doivent jamais migrer vers la table \texttt{profiles}. Toute
vérification d'autorisation côté client (\texttt{hasRole}) est un confort d'UX~;
la véritable barrière de sécurité reste la \textbf{RLS} et les fonctions
\texttt{SECURITY DEFINER} côté base.
\end{infobox}
